\scanpage{29}

\deleted{\unsure{Il est|Il s'|Je dis} \unsure{Comme|On a|que} $x,y\notin$}

Par construction, on a $x\in K$, $x\notin K^p$, \inserted{(cf le NB marginal)} et $y\notin K^p$, $y\in K^p(x)$, donc \deleted{$\alpha,\beta$} \inserted{les $a_i$} sont les \uline{seuls} $\unsure{c|\epsilon|\in|\vphantom{c}}\in K^p$ tels qu'on ait \deleted{$y=\alpha x+\beta$. Evident} \inserted{$y=\sum_0^{p-1}a_i x^i$}

Comme $x\notin K^p$, on a $x\notin K^p\cap k$, mais bien sûr $x^p\in K^p\cap k$, Si donc $(x,y)$ n'\deleted{était} \deleted{\unsure{lié|l'était}} p-libres sur $K^p\cap k$, on aurait une relation $y=\sum_0^{p-1}a'_i x^i$, les $a'_i\in K^p\cap k$, donc \deleted{les} \inserted{on aurait} $a'_i=a_i$, donc $a_i\in k=\mathbb F(x,y)$, absurde.

N.B. Conditions essentielles

\[ \left\{\begin{aligned} &(x,y)\quad\text{p libres sur }\lambda=\mathbb F(a_1,\ldots,a_{p-1})\\ &\deleted{\unsure{a_1=|a_1|a_i}}\ \lambda\not\subset\mathbb F(x,y)\end{aligned}\right. \]

\deleted{Il suffit \inserted{pour} que} exemple \deleted{que $x,y$ et \inserted{$p$ indép. sur} \unsure{$p$ libre sur un|$p$ libre ds un} $\lambda=\mathbb F(a_1,\ldots,a_{p-1})$, et \unsure{par} prendre $\lambda$ \unsure{l'un au seul|l'un au sens}} Prendre $x,y,a_1$ \deleted{alg} p-indép. sur $\mathbb F$, \inserted{dans $\mathbb F(x,y,a)$} et prendre $a_i=0$ pour $2\leq i\leq p-1$. Bien entendu, $\mathbb F$ peut être n'importe quel corps parfait dans \unsure{tt ça.|tt çà.|la sq.}.

On peut aussi prendre $x\unsure{+|\ \text{et}\ |,}y$ alg. indép sur $\mathbb F$, et $a_1$ alg\textsuperscript{\unsure{b|q|\vphantom{b}}} \deleted{séparable} sur $\mathbb F$, $a_1\notin\mathbb F$, Alors $K$ est obtenu \inserted{à partir de $k$,} en extrayant une racine $p$\textsuperscript{ième} (celle de $y-\unsure{a|\alpha}x$) dans $k(b)$, où $b$ est alg. séparable sur $k$ [$b=a^{1/p}$].
% archive sheet number: 28
